\documentclass[oneside,12pt,a4paper]{scrreprt} % ändern auf twoside für Druck-Abgabe

\usepackage[T1]{fontenc}
\usepackage[utf8]{inputenc}
\usepackage[ngerman]{babel}
\usepackage{babelbib}
%% Alternative, APA 7 Citation Style
%%   Dann `babelbib` oben auskommentieren
%%   und weiter unten beim Literaturverzeichnis auch noch ein-/auskommentieren
%% Bitte beachte, dass mit natbib für die Zitation wie Holz (2025) der Befehl \citep verwendet werden muss, siehe auch https://www.overleaf.com/learn/latex/Natbib_citation_styles
% \usepackage[
%  style=apa,
%  backend=biber,
%  sortcites=true,
%  natbib=true,
%  uniquename=false
% ]{biblatex}
% \addbibresource{literatur.bib}  % deine Bib-Datei
\usepackage{parskip}
\usepackage{microtype}
\usepackage{graphicx} % Zum Einbinden von Grafiken
\usepackage[dvipsnames]{xcolor}
\usepackage[colorlinks=true,linkcolor=Black,citecolor=MidnightBlue,urlcolor=MidnightBlue]{hyperref}
\usepackage[all]{hypcap}
\usepackage{pgfplots} \pgfplotsset{compat=1.9}
\usepackage{helvet} % Schönere SansSerif-Schrift
\usepackage{times}  % Schönere Serif-Schrift
\usepackage{enumitem}
\usepackage{blindtext} % sollte am Ende nicht mehr benötigt werden ;)

\usepackage{amsmath}
\usepackage{amssymb}

\pagestyle{headings}

\graphicspath{ {figures/} } % Pfad-Prefix für einzubindende Grafiken. Es sind auch mehrere Pfade möglich, diese müssen jeweils in eigenen {Klammern} stehen.

\setkomafont{disposition}{\normalcolor\bfseries} % überall Serifen verwenden
% oder
%\renewcommand{\familydefault}{\sfdefault} % überall Sans-Serif verwenden

% PDF-Optionen (werden in den Dateieigenschaften angezeigt)
\hypersetup{
pdftitle={Titel der Arbeit},
pdfauthor={Vor- und Nachname},
pdfsubject={Bachelorarbeit Informatik},
pdfpagelayout=TwoColumnRight
}

%%% Eigene Makros
\newcommand{\qq}[1]{\glqq #1\grqq} % \qq{Text in Anführungszeichen}

\begin{document}

%%% Titelseite
\begin{titlepage}
\begin{center}
\LARGE Pädagogische Hochschule Ludwigsburg \\
\large Fakultät II: Kultur- und Naturwissenschaften \\
Institut für Informatik\\
[1cm]

\Large Bachelor-/Masterstudiengang ---~\small{(bitte Studiengang angeben)}
\\
\Large Bachelor of Arts/Master of Education/Master of Arts
\\[1cm]
\huge Bachelor-/Masterarbeit Informatik
\large\\[0.75cm]
\rule{0.5\textwidth}{1pt}
\\[0.75cm]
%[1.5cm]
\large
Genehmigungsdatum
\\[0.5cm]
\Large\textbf{Thema/Titel der Arbeit (der geht wohl in den\\ meisten Fällen über mehr als eine Zeile)}
\\\small\textrm(genauer Wortlaut siehe Genehmigung; Groß- und Kleinschreibung beachten, am Ende kein Punkt)\\
[0.5cm]
\large
Abgabedatum
\\[1cm]
\small
\textbf{Prüfer*in/nen}\\[0.3cm] 
\large Name Prüfer*in 
\\\small (mit genauem Titel und Berufsbezeichung, z.B. Jun.-Prof. Dr. Heiko Holz)
\\[.75cm]
\rule{0.5\textwidth}{1pt}
\\[0.7  5cm]
\small
\textbf{Absolvent}\\[0.3cm]
\large Name, Vorname\\
Matrikelnummer\\
E-Mail-Adresse \small(NameVorname@stud.ph-ludwigsburg.de)\\\textcolor{red}{Private E-Mail ist nicht zulässig!}
%\\[1cm]
\vfill
\fbox{%
\begin{minipage}{\textwidth}
\centering
\large
Einverständnis für Freigabe der Arbeit:\\
$\square$ ja \quad $\square$ nein (da die Arbeit personenbezogene Daten enthält)
\end{minipage}
}

\end{center}
\end{titlepage}

%%% Titelrückseite: Bibliographische Angaben
\thispagestyle{empty}
\vspace*{\fill}
\textbf{Nachname, Vorname:}\\
\emph{Titel der Arbeit}\\
Bachelor-/Masterarbeit Informatik\\
Pädagogische Hochschule Ludwigsburg\\
Bearbeitungszeitraum: Anfangs- -- Enddatum
\newpage

%%% Zusammenfassung (Abstract), hier aus externer Datei eingebunden
% !TEX root = ../ausarbeitung.tex

\begin{abstract}
\section*{Zusammenfassung}
Wissenschaftliche Arbeiten fangen normalerweise mit einer kurzen Zusammenfassung an. Deshalb sollte Ihre Arbeit ebenfalls eine solche Zusammenfassung enthalten. Die Zusammenfassung hat einen ähnlichen Inhalt wie die Motivation, nur viel kürzer. Sie soll kurz beschreiben
\begin{itemize}
\item worum es in der Arbeit geht (was war das zu lösende Problem?),
\item welche Methoden zur Problemlösung angewendet wurden,
\item wie das ganze evaluiert wurde,
\item evtl. welches Ergebnis/ Schlussfolgerungen sich daraus ergeben.
\end{itemize}

\hfil\rule{0.4\textwidth}{0.4pt}

Dieses Dokument soll als Ausgangs-Template für Bachelorarbeiten dienen. Gleichzeitig soll es zeigen, wie so ein \qq{fertiges Dokument} aussehen könnte. Um die Seiten gefüllt zu bekommen, wurde Blindtext verwendet.

Die kurzen Beschreibungen zu den Abschnitten (jeweils über dem Querstrich) wurden \cite{heikoholz2024} entnommen, welches wiederum eine Überarbeitung von~\cite{alexandrakirsch2016} ist. Diese \qq{Hinweise zum Erstellen von Bachelor-/Masterarbeiten im Institut für Informatik der Pädagogischen Hochschule Ludwigsburg} sind aber auch darüber hinaus zu empfehlen.
\end{abstract}
\newpage

%%% Inhaltsverzeichnis
\KOMAoption{toc}{listof,bib} % Abbildungs-/Tabellenverzeichnis, Literaturverzeichnis aufnehmen
\tableofcontents\label{toc}
\cleardoublepage

%%% Hauptteil (mit \input{dateiname} wird die Datei 'dateiname' eingebunden)
% !TEX root = ../ausarbeitung.tex

\chapter{Einleitung}
Die Einleitung ist (ebenso wie die Zusammenfassung) das Aushängeschild Ihrer Arbeit. Nach der Einleitung entscheidet ein Leser, ob er den Rest der Arbeit überhaupt lesen möchte! Daher sollten Sie vor allem die Einleitung mehrfach überarbeiten, damit sie sich angenehm liest und kurz und präzise die folgenden Punkte beschreibt:
\begin{itemize}
  \item Was wurde in der Arbeit untersucht? Was ist daran neu?
  \item evtl: Im Rahmen welches größeren Projektes am Lehrstuhl wurde die Arbeit durchgeführt?
  \item Warum ist die Arbeit wichtig?
  \item Warum ist die Arbeit schwierig?
  \item Was sind die (wissenschaftlichen) Beiträge der Arbeit?
  \item Welche ähnlichen Arbeiten/Lösungsansätze für ähnliche Probleme gibt es schon? Wie unterscheiden sich diese von dem in der Arbeit gewählten Lösungsweg? (Diese Fragen können auch am Ende der Arbeit oder in einem eigenen Gliederungspunkt beantwortet werden.)
\end{itemize}

Der letzte Abschnitt der Einleitung enthält oft einen Überblick über die darauffolgenden Kapitel, als eine Art \qq{Karte} für den Leser. Dieser Teil ist aber eher eine \qq{Pflichtübung}, die weniger wichtig ist als die obigen Inhalte.

Die gesamte Ausarbeitung sollte eine schlüssige \qq{Geschichte} erzählen und die Einleitung gibt diese Geschichte quasi schon vor:
\begin{itemize}
  \item Was wurde in der Arbeit untersucht?\\
  Wie passt die Arbeit ins Gesamtbild der Wissenschaftslandschaft? Diese Frage sollte am Ende bei der Diskussion/Zusammenfassung noch einmal aufgegriffen werden.
  \item Warum ist die Arbeit wichtig?\\
  Wenn sie nicht wichtig wäre, müsste man sie nicht machen. Die konkrete Fragestellung sollte sich idealerweise aus dem aktuellen Stand der Wissenschaft ergeben ($\rightarrow$ related work)
  \item Warum ist die Arbeit schwierig?\\
  Wenn sie einfach wäre, wäre sie nicht interessant. Welche Herausforderungen sind also konkret in Angriff zu nehmen? \emph{Wie} Sie die Herausforderungen gelöst haben, beschreiben Sie im Hauptteil der Arbeit.
  \item Was sind die (wissenschaftlichen) Beiträge der Arbeit?\\
  Man kann in der Einleitung auch beschreiben, was die Ziele der Arbeit sind und die Beiträge in der Zusammenfassung aufschreiben. Die Beiträge der Arbeit sollten genau den gesetzten Zielen entsprechen.\\
  Die Ziele, die Sie sich in der Einleitung setzen, führen zu entsprechenden Evaluationsmethoden. Z.B.\ wenn ein Ziel/Beitrag ist, dass die von Ihnen entwickelte Nutzeroberfläche barrierefrei ist, dann müssen Sie in der Evaluation zeigen, dass sie dies tatsächlich ist (bzw.\ zu welchem Grad sie es ist). Wissenschaftliche Beiträge sind Dinge, die Sie selbst gemacht haben und die die Wissenschaft ein Stückchen weiter voranbringen. Dies kann eine ausgiebige Literaturrecherche sein, ein Vergleich von Systemen/Nutzeroberflächen, die Implementierung eines Systems, ein Nutzertest, eine Umfrage, etc. Ihre Arbeit sollte 2--4 (3 ist perfekt) Beiträge enthalten.
\end{itemize}
\cleardoublepage

% !TEX root = ../ausarbeitung.tex

\chapter{Stand der Forschung}
Hier zeigen Sie, dass Sie über Ihr Themengebiet gut informiert sind. Sie können entweder den Stand der Forschung dafür heranziehen, um Ihr Thema zu rechtfertigen (\qq{Warum ist es wichtig?}) oder Sie können die Literatur als Grundlage Ihrer Diskussion verwenden (Wie ordnen sich Ihre Beiträge in die Wissenschaftslandschaft allgemein ein?), eine Mischung ist auch möglich.

Dazu gehören natürlich Referenzen zu anderen Werken. Für deren Verwaltung empfiehlt sich BibTeX oder BibLaTex, die entsprechende Datei \verb|literatur.bib| (kann natürlich auch anders benannt sein) ist in der Vorlage schon enthalten. Zur Einhaltung der Syntax bietet sich ein Online-Editor wie \cite{BibTexOnlineEditor} an; TeXstudio \cite{texstudio} hält im Menü auch Hilfe beim Erstellen der Bibliographie-Einträge bereit. Für Bücher bietet z.B. auch \url{https://books.google.de/} fertige BibTeX-Einträge an.

Die Einbindung in den Text erfolgt dann mit \verb|\cite{nielsen1994usability}| (bzw. \verb|\citep{nielsen1994usability}| für APA-Zitation mit BiBLaTex und natbib, siehe \url{https://www.overleaf.com/learn/latex/Natbib_citation_styles}), wobei der Text in den Klammern durch das in der \verb|.bib|-Datei vergebene Kürzel ersetzt werden müssen. Das Ergebnis ist dann eine Referenz zum (zumindest im Bereich \qq{Usability Engineering}) fast unverzichtbaren gleichnamigen Buch \cite{nielsen1994usability} von Jakob Nielsen.

Für die Literaturverwaltung empfehlen wir \href{https://www.zotero.org/}{Zotero}~\cite{zotero}. In Zotero können Referenzen auch ganz einfach als BibTex oder BibLaTex exportiert werden. Hierfür empfehlen wir das Plugin \href{https://github.com/retorquere/zotero-better-bibtex}{Better BibTex for Zotero}. Am Institut für Informatik stellen wir bereits einige Literatur für Abschlussarbeiten in der Zotero-Gruppe \href{https://www.zotero.org/groups/5530859/}{
Abschlussarbeiten bei und Projekte von Heiko Holz} bereit. Sie können, nachdem Sie sich bei Zotero registriert haben, den Beitritt in die Gruppe beantragen.

Achtung: Ihre Ausarbeitung sollte sich (im Gegensatz zu diesem Template) weniger auf Internet-Quellen, als auf Bücher und Paper stützen!
\\

\hfil\rule{0.4\textwidth}{0.4pt}

\Blindtext[5]

\cleardoublepage

% !TEX root = ../ausarbeitung.tex

\chapter{Herangehensweise}

Im Hauptteil beschreiben Sie Ihre praktische Arbeit. Code gehört normalerweise nicht in eine Ausarbeitung. Ausnahmen sind Algorithmen, die für Sie wichtig waren (dann in möglichst übersichtlichem Pseudo-Code). Kleine Code-Stücke können auch zur Illustration oder als Beispiel eingebaut werden. Längere Code-Stücke können im Anhang untergebracht werden. Sie sollten jedoch nicht den gesamten Code im Anhang abdrucken. Ihren Code geben Sie bitte dennoch mit ab, am besten auf einer CD.

Der Detailgrad sollte so sein, dass ein Leser die gleiche Arbeit noch einmal nachimplementieren könnte. Insbesondere sollten alle Parameter, von denen die Funktionsweise des Systems abhängt, explizit genannt sein. Bei der Evaluation muss bei jedem Versuch angegeben werden, mit welchen Parametern gearbeitet wird.


\hfil\rule{0.4\textwidth}{0.4pt}

\section{Einführung in \LaTeX}
Wenn Sie dieses Template gefunden haben, werden Sie schon erkannt haben, dass \LaTeX\ nicht ganz unwichtig ist. Falls Sie noch kaum oder keine Erfahrung damit haben, bietet sich der Kurs \qq{Einführung in \LaTeX} von Thorsten Nagel an. Er findet meistens mehrmals im Semester statt, die genauen Termine entnehmen Sie bitte dem Vorlesungsverzeichnis. Das Skript zu diesem Kurs ist online erhältlich \cite{thorstennagel2015} und enthält noch viele weitere Informationen, die in der folgenden Kurzzusammenfassung nicht berücksichtigt werden konnten.

\section{Sinnvolle Programme}
Falls Sie \verb|.tex|-Dateien nicht unbedingt in einem normalen Texteditor schreiben und auf der Kommandozeile kompilieren möchten, bietet sich ein spezieller \LaTeX-Editor wie z.B. TeXstudio \cite{texstudio} an.

\section{Wichtige \LaTeX-Befehle}
\subsection{Textgliederung}
Mit dem Befehl \verb|\section{Überschrift}| lässt sich ein neuer Abschnitt beginnen, der automatisch nummeriert und im \hyperref[toc]{Inhaltsverzeichnis} eingetragen wird.
Gleiches gilt für einen Unterabschnitt, der mit \verb|\subsection{Unterüberschrift}| begonnen wird.

\LaTeX\ bietet noch weitere Gliederungsebenen bis hin zum Unterabsatz, die aber standardmäßig nicht mehr nummeriert werden und standardmäßig auch nicht im Inhaltsverzeichnis auftauchen. Im Normalfall sind diese aber nicht unbedingt nötig.
\subsubsection{Unterunterüberschrift}
\paragraph{Absatztitel} Absatztext
\subparagraph{Unterabsatz} Absatztext

Um in der PDF-Datei einen Zeilenumbruch zu erhalten, schreiben Sie im Quelltext \verb|\\|, für einen neuen Absatz lassen Sie eine Zeile frei. Ein einzelner Zeilenumbruch im Quelltext hat keine Auswirkungen.

\section{Formatierungen, Listen, Aufzählungen}
Um die \LaTeX-Befehle für die folgenden Formatierungen zu sehen, schauen Sie sich einfach den Quellcode an.

Text kann man \textbf{fett} oder \textit{kursiv} schreiben. Sie können Text auch explizit \emph{betonen}, was dann ebenfalls kursiv dargestellt wird.

Listen sind natürlich auch möglich, diese können beliebig verschachtelt werden.
\begin{itemize}
\item Listenpunkt 1
\item Listenpunkt 2
\begin{itemize}
\item Unterpunkt
\end{itemize}
\end{itemize}

Oder eine nummerierte Aufzählung:
\begin{enumerate}
\item Listenpunkt 1
\item Listenpunkt 2
\begin{enumerate}
\item Unterpunkt
\end{enumerate}
\end{enumerate}

\section{Tabellen}
Tabellen können wie folgt erstellt werden (siehe \autoref{tab:tabelle-1}):
\begin{table}[htb]
\centering
\caption{Formatierungsparameter zur Spaltenausrichtung in Tabellen\label{tab:tabelle-1}}
\begin{tabular}{l|p{3.5cm}lcr}
	                     & \textbf{Spalte 1}                   & \textbf{Spalte 2} & \textbf{Spalte 3} & \textbf{Spalte 4} \\ \hline
	\textbf{Parameter}   & \verb|p{3.5cm}|                     & \verb|l|          &     \verb|c|      &          \verb|r| \\
	\textbf{Ausrichtung} & linksbündig mit vorgegebener Breite & linksbündig       &     zentriert     &      rechtsbündig
\end{tabular}
\end{table}

Für \qq{schönere} Tabellen empfiehlt sich das \LaTeX-Paket \verb|booktabs|, das in der Präambel mit \verb|\usepackage{booktabs}| eingebunden werden kann. Insbesondere die zugehörige Anleitung \cite{booktabs} ist lesenswert, da sie recht einfach vermittelt, was eine \qq{schöne} Tabelle ausmacht.

\section{Mathematische Gleichungen}
Gleichungen oder sonstige Formeln im Fließtext, wie z.B. $f(x)=ax+b$, können Sie einfach in Dollarzeichen \verb|$...$| einfassen. Größere Formeln können auch in einer eigenen Zeile stehen:
\begin{displaymath}
g(x) = \frac{1}{2\pi\sigma}\cdot e^{-\frac{(x-\mu)^2}{2\sigma^2}}
\end{displaymath}
Falls Sie in Ihrer Ausarbeitung viel Mathematik benötigen, lohnt sich auch ein Blick auf das \LaTeX-Paket \verb|amsmath|.
\section{Abbildungen}
Abbildungen und Diagramme können Sie direkt in einer \verb|.tex|-Datei \qq{malen}. Dazu bietet TikZ viele Möglichkeiten, eine Grafik mit Text zu beschreiben. Falls Sie nicht jeden Strich in einem Diagramm selbst definieren wollen, bietet sich die Verwendung von Pgfplots in Form der \verb|axis|-Umgebung an, wie in \autoref{fig:diagramm} gezeigt.

Falls Sie sich näher mit der Materie beschäftigen wollen, bietet sich das Manual für TikZ und PGF\cite{pgfmanual} an. Dieses Dokument ist zwar sehr umfangreich (über 1000 Seiten), bietet aber viele gut erklärte Beispiele, auch und gerade für den Einstieg.

Die folgenden Abbildungen wurden in jeweils eine \verb|figure|-Umgebung gesteckt. Diese erzeugt eine \qq{floatende} Abbildung, die von \LaTeX\ automatisch an einer \qq{geeigneten} Stelle platziert wird, also nicht zwangsläufig dort, wo sie im Quelltext definiert wurde. Wenn Sie der Abbildung eine \verb|\caption{...}| und ein \verb|\label{key}| verpassen, können Sie über dieses Label im Text per \verb|\ref{key}| oder \verb|\autoref{key}| darauf verweisen. Bei ersterem wird nur die Nummer (\qq{\ref{fig:diagramm}}), bei letzterem zusätzlich noch der Typ (\qq{\autoref{fig:diagramm}}) ausgegeben. Analog zur \verb|figure|-Umgebung wurde oben für die Tabelle die \verb|table|-Umgebung verwendet.
\begin{figure}[htb]
\centering
\begin{tikzpicture}
  \begin{axis}[xlabel=$x$,ylabel=$y$,width=0.75\textwidth]
    \addplot[smooth,mark=*,blue] plot coordinates {
        (0,2)
        (2,3)
        (3,1)
    };
    \addlegendentry{Case 1}

    \addplot[smooth,color=red,mark=x] plot coordinates {
            (0,0)
            (1,1)
            (2,1)
            (3,2)
        };
    \addlegendentry{Case 2}
  \end{axis}
\end{tikzpicture}
\caption{Diagramm aus den TikZ-Beispielen für Pgfplots\label{fig:diagramm}}
\end{figure}

Falls Sie eine Abbildung schon als separate Datei (z.B. im PDF-Format) vorliegen haben, können Sie diese mit \verb|\includegraphics{}| und dem Dateinamen einbinden, wie in \autoref{fig:smileys} geschehen. Dabei wird in den in der Präambel mit \verb|\graphicspath{...}| vorgegebenen Pfaden gesucht. Die Dateiendung wird automatisch \qq{erraten}, wenn sie nicht angegeben wird.
\begin{figure}[htb]
	\centering
	\includegraphics[width=0.75\textwidth]{smileys}
	\caption{Darstellung des Wohlbefindens anhand von Smileys\label{fig:smileys}}
\end{figure}

\cleardoublepage

% !TEX root = ../ausarbeitung.tex

\chapter{Evaluation}
Der Evaluationsteil sollte zunächst die Forschungsfragen/anvisierten Beiträge aus der Einleitung aufgreifen und daraus die verwendeten Evaluationsmethoden ableiten. Z.B. \qq{Das Ziel war es eine barrierefreie, effizente Bedienoberfläche zu erstellen, die auf verschiedenen Endgeräten läuft.} Daraus ergeben sich folgende Evaluationsziele und Methoden:
\begin{itemize}
  \item barrierefrei $\rightarrow$ Heuristiken zur Barrierefreiheit nach W3C oder BITV
  \item effizient $\rightarrow$ Messung der Antwortzeit für verschiedene Anfragen
  \item verschiedene Endgeräte $\rightarrow$ Anzeigebeispiele von verschiedenen Endgeräten oder Test durch entsprechendes Tool, das verschiedene Endgeräte simuliert.
\end{itemize}

Es bietet sich an, die Durchführung der Evaluation und ihre Beschreibung parallel zu bearbeiten, denn oft fällt erst beim Schreiben auf, dass die Evaluationsmethode nicht zu den Zielen der Arbeit passt. Die Beschreibung sollte alle Parameter enthalten, die den Versuch ausmachen, d.h. Programmparameter, Testbedingungen, Anzahl Testläufe etc.

\hfil\rule{0.4\textwidth}{0.4pt}

\Blindtext[6]
\cleardoublepage

% !TEX root = ../ausarbeitung.tex

\chapter{Diskussion}
Die Diskussion kann als Teil des Evaluations- oder Schlusskapitels oder als eigenständiges Kapitel aufgeführt werden. Wichtig ist, dass Sie Ihre Evaluationsergebnisse realistisch einschätzen und ins Verhältnis zum Stand der Technik setzen. Achten Sie besonders darauf, aus den Daten Ihrer Evaluation keine Wunschergebnisse abzulesen, die nicht in den Daten sind (wenn Ihre Testnutzer Ihr neuimplementiertes System nicht besser bedienen können als ein vorhandenes, dann ist das eben so). Gerade unerwartete bzw. \qq{negative} Ergebnisse (z.B. das neue System ist nicht besser als vorhandene) bringen wissenschaftliche Erkenntnisse: man stellt damit fest, dass der gewählte Weg nicht zum gewünschten Ergebnis führt und man generiert damit neue Fragen, z.B. warum der Weg nicht funktioniert hat, obwohl er vor dem Test als überlegen erachtet wurde.

Es kann auch sein, dass verschiedene Evaluationsformen Unterschiede offenbaren. Z.B. kann es sein, dass die Nutzbarkeit des implementierten Systems nicht besser ist als bei anderen Ansätzen, aber dass es deutlich einfacher zu warten ist.


Im letzten Teil runden Sie Ihre Arbeit ab, in dem Sie Ihre Argumentation aus der Einleitung aufgreifen und mit konkreten Daten aus Ihrem Hauptteil und der Evaluation untermauern. Auch hier können Sie Bezug zur Literatur nehmen. Am Ende sollten Sie einen Ausblick über weitere Forschungsthemen geben. Dabei aufpassen, dass es nicht so klingt wie \qq{mir ist die Zeit ausgegangen und folgendes habe ich nicht mehr geschafft}. Eine gute wissenschaftliche Arbeit wirft mehr Fragen auf als sie beantwortet. Es sollte also eher klingen nach \qq{meine Arbeit hat \dots gezeigt. Daraus ergeben sich weitere interessante Fragen \dots}.
\cleardoublepage

%%% Literaturverzeichnis, lädt die Datei literatur.bib
\bibliographystyle{babplain} % "babplain" benötigt das Paket babelbib
\bibliography{literatur}
%% Alternative, APA 7 Citation Style
%%   Dann obige beide Zeilen auskommentieren
%\printbibliography

\cleardoublepage

%%% Selbständigkeitserklärung
\thispagestyle{empty}
\section*{Erklärung}
An dieser Stelle ist die gültige Fassung der \href{https://media.ph-ludwigsburg.de/downloadzentrum/objekte/1410173-Eigenstndigkeitserklrung-fr-schriftliche-Prfungen
}{Eigenständigkeitserklärung für schriftliche Prüfungen} der Pädagogischen Hochschule Ludwigsburg einzubinden.
\hrule

\noindent

%\vspace{1cm}

\noindent
Name, Vorname:

\noindent
Matrikelnummer: 
\vspace{0.1cm}
\hrule
\vspace{0.5cm}

\noindent
\begin{itemize}[label=$\square$, leftmargin=*, nosep]
    \item Bachelorarbeit
    \item Masterarbeit
    \item schriftlichen Modulprüfung (Hausarbeit, Portfolio, etc.)
    \item meinen Anteil der Gruppenarbeit als schriftliche Modulprüfung\\
    (der persönliche Teil muss gekennzeichnet sein)
\end{itemize}

\vspace{0.5cm}

\noindent
Hiermit versichere ich, dass ich die vorliegende Arbeit mit dem Titel:

\vspace{0.5cm}

\noindent
\hrule
\noindent
selbstständig und ohne unzulässige fremde Hilfe oder Hilfsmittel verfasst habe.

\noindent
Ich habe keine anderen als die angegebenen Quellen und Hilfsmittel benutzt.

\noindent
Alle verwendeten Quellen und sämtliche Literatur habe ich vollständig im Quellen- bzw. Literaturverzeichnis aufgeführt. Ich habe alle wörtlichen Zitate und sinngemäßen Verweise gekennzeichnet.

\noindent
Ich habe alle Hilfsmittel aufgeführt, die ich genutzt habe (hierzu zählen auch generative KI oder Dialogsysteme). [Hilfsmittel sind im Anschluss an das Quellen- bzw. Literaturverzeichnis aufzuführen.]

\noindent
Diese Versicherung gilt für sämtliche Bestandteile meiner Arbeit.

\noindent
Meine Qualifikationsarbeit hat in gleicher oder ähnlicher Form noch keiner Prüfungsbehörde an dieser oder einer anderen Hochschule vorgelegen. Ebenfalls erkläre ich, dass ich noch keine Abschlussprüfung in einem gleichen oder vergleichbaren Studiengang endgültig nicht bestanden habe und mich nicht in einem laufenden Prüfungsverfahren in einem gleichen oder vergleichbaren Studiengang befinde.

\vspace{2cm}

\noindent
\begin{tabular}{@{}p{5cm}p{5cm}@{}}
\hrulefill & \hrulefill \\
Ort, Datum & Unterschrift
\end{tabular}
\vfill
Bearbeitungsstand 01.04.2024

\end{document}